
\documentclass{article}
\usepackage{colortbl}
\usepackage{makecell}
\usepackage{multirow}
\usepackage{supertabular}

\begin{document}

\newcounter{utterance}

\centering \large Interaction Transcript for game `cladder', experiment `full\_v1.5\_default', episode 9771 with refineonly{-}Qwen3.5{-}27B.
\vspace{24pt}

{ \footnotesize  \setcounter{utterance}{1}
\setlength{\tabcolsep}{0pt}
\begin{supertabular}{c@{$\;$}|p{.15\linewidth}@{}p{.15\linewidth}p{.15\linewidth}p{.15\linewidth}p{.15\linewidth}p{.15\linewidth}}
    \# & \multicolumn{2}{c}{Player} && \multicolumn{2}{c}{Game Master} \\
    \hline

    \theutterance \stepcounter{utterance}  
    & & & \multicolumn{4}{p{0.6\linewidth}}{
        \cellcolor[rgb]{0.9,0.9,0.9}{
            \makecell[{{p{\linewidth}}}]{
                \texttt{\tiny{[P1$\langle$GM]}}
                \texttt{You are an expert in causal inference. The following question is not a typical commonsense query, but rather a meticulously designed question created by a professor specializing in causal inference, intended to assess the students' mastery of the course content.} \\
\texttt{Imagine a self{-}contained, hypothetical world with only the following conditions, and without any unmentioned factors or causal relationships: Jyka has a direct effect on yupt and kwox. Yupt has a direct effect on kwox. We know that jyka causes yupt. jyka or yupt causes kwox. Would an individual is kwox if not jyka instead of jyka?} \\
\texttt{Start your answer with 'Yes' or 'No', followed by additional reasoning or evidence to support your explanation.} \\
            }
        }
    }
     \\ \\

    \theutterance \stepcounter{utterance}  
    & & \multicolumn{4}{p{0.6\linewidth}}{
        \cellcolor[rgb]{0.95,0.95,0.95}{
            \makecell[{{p{\linewidth}}}]{
                \texttt{\tiny{[P1$\leftrightarrow$P]}}
                \texttt{{[}DRAFT{]}} \\
\texttt{No} \\
\\ 
\texttt{**Reasoning:**} \\
\\ 
\texttt{To determine whether an individual would be "kwox" if they were "not jyka" instead of "jyka," we must analyze the causal structure and the specific conditions provided in this hypothetical world.} \\
\\ 
\texttt{1.  **Causal Structure**:} \\
\texttt{    *   \$Jyka \textbackslash{}rightarrow Yupt\$ (Jyka causes Yupt).} \\
\texttt{    *   \$Jyka \textbackslash{}rightarrow Kwox\$ (Jyka has a direct effect on Kwox).} \\
\texttt{    *   \$Yupt \textbackslash{}rightarrow Kwox\$ (Yupt has a direct effect on Kwox).} \\
\texttt{    *   The prompt states: "Jyka or Yupt causes Kwox." In the context of causal graphs, this usually implies that Kwox is the outcome of these parents.} \\
\\ 
\texttt{2.  **The Counterfactual Scenario**:} \\
\texttt{    *   **Actual World**: The individual is "Jyka". Consequently, they are "Yupt" (due to \$Jyka \textbackslash{}rightarrow Yupt\$) and "Kwox" (due to direct effect from Jyka and/or indirect effect via Yupt).} \\
\texttt{    *   **Counterfactual World**: We intervene to set \$Jyka = \textbackslash{}text\{False\}\$ (not Jyka).} \\
\\ 
\texttt{3.  **Analyzing the Effect of the Intervention**:} \\
\texttt{    *   Since \$Jyka\$ causes \$Yupt\$, if we set \$Jyka = \textbackslash{}text\{False\}\$, the value of \$Yupt\$ changes. In a deterministic causal model (implied by the binary nature of the question "is kwox"), if the cause (\$Jyka\$) is removed, the effect (\$Yupt\$) typically ceases to exist or becomes False, unless there is another unmentioned cause for \$Yupt\$. The prompt explicitly states "without any unmentioned factors or causal relationships." Therefore, if \$Jyka\$ is False, \$Yupt\$ becomes False.} \\
\texttt{    *   Now consider \$Kwox\$. \$Kwox\$ depends on \$Jyka\$ and \$Yupt\$.} \\
\texttt{        *   The direct path \$Jyka \textbackslash{}rightarrow Kwox\$ is broken because \$Jyka\$ is now False.} \\
\texttt{        *   The indirect path \$Jyka \textbackslash{}rightarrow Yupt \textbackslash{}rightarrow Kwox\$ is also broken because \$Yupt\$ is now False (as it relies entirely on \$Jyka\$).} \\
\texttt{    *   The prompt states "Jyka or Yupt causes Kwox." This implies that for \$Kwox\$ to be true, at least one of its causes (\$Jyka\$ or \$Yupt\$) must be present.} \\
\texttt{    *   In the counterfactual world where \$Jyka\$ is False, both \$Jyka\$ and \$Yupt\$ are False.} \\
\\ 
\texttt{4.  **Conclusion**:} \\
\texttt{    *   Since both necessary causes (\$Jyka\$} \\
            }
        }
    }
    & \\ \\

    \theutterance \stepcounter{utterance}  
    & \multicolumn{4}{p{0.6\linewidth}}{
        \cellcolor[rgb]{0.561,0.737,0.561}{
            \makecell[{{p{\linewidth}}}]{
                \texttt{\tiny{[P1$\rangle$GM]}}
                \texttt{No} \\
\\ 
\texttt{**Reasoning:**} \\
\\ 
\texttt{To determine whether an individual would be "kwox" if they were "not jyka" instead of "jyka," we must analyze the causal structure and the specific conditions provided in this hypothetical world.} \\
\\ 
\texttt{1.  **Causal Structure**:} \\
\texttt{    *   \$Jyka \textbackslash{}rightarrow Yupt\$ (Jyka causes Yupt).} \\
\texttt{    *   \$Jyka \textbackslash{}rightarrow Kwox\$ (Jyka has a direct effect on Kwox).} \\
\texttt{    *   \$Yupt \textbackslash{}rightarrow Kwox\$ (Yupt has a direct effect on Kwox).} \\
\texttt{    *   The prompt states: "Jyka or Yupt causes Kwox." In the context of causal graphs, this implies that Kwox is the outcome of these parents.} \\
\\ 
\texttt{2.  **The Counterfactual Scenario**:} \\
\texttt{    *   **Actual World**: The individual is "Jyka". Consequently, they are "Yupt" (due to \$Jyka \textbackslash{}rightarrow Yupt\$) and "Kwox" (due to direct effect from Jyka and/or indirect effect via Yupt).} \\
\texttt{    *   **Counterfactual World**: We intervene to set \$Jyka = \textbackslash{}text\{False\}\$ (not Jyka).} \\
\\ 
\texttt{3.  **Analyzing the Effect of the Intervention**:} \\
\texttt{    *   Since \$Jyka\$ causes \$Yupt\$, if we set \$Jyka = \textbackslash{}text\{False\}\$, the value of \$Yupt\$ changes. In a deterministic causal model (implied by the binary nature of the question "is kwox"), if the cause (\$Jyka\$) is removed, the effect (\$Yupt\$) typically ceases to exist or becomes False, unless there is another unmentioned cause for \$Yupt\$. The prompt explicitly states "without any unmentioned factors or causal relationships." Therefore, if \$Jyka\$ is False, \$Yupt\$ becomes False.} \\
\texttt{    *   Now consider \$Kwox\$. \$Kwox\$ depends on \$Jyka\$ and \$Yupt\$.} \\
\texttt{        *   The direct path \$Jyka \textbackslash{}rightarrow Kwox\$ is broken because \$Jyka\$ is now False.} \\
\texttt{        *   The indirect path \$Jyka \textbackslash{}rightarrow Yupt \textbackslash{}rightarrow Kwox\$ is also broken because \$Yupt\$ is now False (as it relies entirely on \$Jyka\$).} \\
\texttt{    *   The prompt states "Jyka or Yupt causes Kwox." This implies that for \$Kwox\$ to be true, at least one of its causes (\$Jyka\$ or \$Yupt\$) must be present.} \\
\texttt{    *   In the counterfactual world where \$Jyka\$ is False, both \$Jyka\$ and \$Yupt\$ are False.} \\
\\ 
\texttt{4.  **Conclusion**:} \\
\texttt{    *   Since both necessary causes (\$Jyka\$ and} \\
            }
        }
    }
    & & \\ \\

    \theutterance \stepcounter{utterance}  
    & & & \multicolumn{2}{p{0.3\linewidth}}{
        \cellcolor[rgb]{0.9,0.9,0.9}{
            \makecell[{{p{\linewidth}}}]{
                \texttt{\tiny{[GM$|$GM]}}
                \texttt{No} \\
            }
        }
    }
    & & \\ \\

    \theutterance \stepcounter{utterance}  
    & & & \multicolumn{2}{p{0.3\linewidth}}{
        \cellcolor[rgb]{0.9,0.9,0.9}{
            \makecell[{{p{\linewidth}}}]{
                \texttt{\tiny{[GM$|$GM]}}
                \texttt{no} \\
            }
        }
    }
    & & \\ \\

    \theutterance \stepcounter{utterance}  
    & & & \multicolumn{2}{p{0.3\linewidth}}{
        \cellcolor[rgb]{0.9,0.9,0.9}{
            \makecell[{{p{\linewidth}}}]{
                \texttt{\tiny{[GM$|$GM]}}
                \texttt{game\_result = WIN} \\
            }
        }
    }
    & & \\ \\

\end{supertabular}
}

\end{document}
